RDB保存某个时点的状态，AOF保存用于重建状态的修改记录。命令执行后，Redis按传播规则把修改编码到AOF缓冲区，之后写入文件，并根据策略安排同步。执行成功、写入用户态缓冲、操作系统接受文件写入以及存储同步完成，分别处在这条路径上。AOF使快照之后的修改有持久记录，实际可恢复前缀由同步策略和故障范围决定~\cite{redis8102-persistence,redis8102-source}。

\subsection{三种同步策略与响应位置}

配置\code{appendfsync always}在正常路径下把相应AOF写入同步与客户端回复发送的顺序连接起来。事件循环在发送待回复数据前处理AOF，正常的\code{always}路径要求相应同步完成后再推进回复发送。应用观察到成功回复时，可以据此判断本地相应同步层已经越过该位置，存储设备失效、虚报同步或整台设备丢失属于另外的故障前提。

\code{appendfsync everysec}通过后台同步减小逐次等待成本。主线程把数据写入文件后可以继续处理请求，同步工作随后推进。通常所说的约一秒损失窗口是这一策略的运行目标，后台同步迟缓、文件写入被推迟或存储异常时，实际状态要结合AOF缓冲、同步进度和错误判断。当前实现会考虑正在进行的后台同步而推迟写入，并在等待过久时继续写入，配置名称并没有给出严格的一秒上界。

\code{appendfsync no}把同步时机主要交给操作系统。进程被杀死以后，内核页缓存可能保留，重启进程看见全部记录很正常。机器断电则会丢失另一层尚未落到可靠介质的内容。一次进程退出后数据还在，只能说明该次进程故障下的恢复结果。

\CEighteenFigure{ch18-08}{AOF缓冲、文件写入、同步与客户端观察}{fig:18-aof-sync}

\begin{table}[htbp]
\centering\small
\caption{AOF同步策略与恢复位置}\label{tab:18-aof-policy}
\begin{tabularx}{\linewidth}{@{}p{.14\linewidth}XX@{}}
\toprule
同步策略 & 正常路径中的同步安排 & 购买状态的恢复判断\\
\midrule
{\ttfamily always} & 频繁同步并约束相应回复发送 & 核对本地同步成功及所采用的存储故障模型\\
{\ttfamily everysec} & 文件写入与后台周期同步分开推进 & 已返回写入可能位于最后可靠同步位置之后\\
{\ttfamily no} & Redis不按每次或每秒策略主动等待同步 & 根据操作系统与设备的实际持久前缀恢复\\
\bottomrule
\end{tabularx}
\end{table}

配置例外也进入同一个判断。例如\code{no-appendfsync-on-rewrite}会改变后台重写期间的同步安排。磁盘满、短写和同步失败各有错误路径，服务器会记录错误并依据策略拒绝、停止或退出相应路径。AOF写入状态、重写进度、同步延迟与磁盘余量共同说明最近修改能否进入可靠前缀，恢复者据这个前缀确定缺失范围。

\subsection{重写把长历史压缩成当前状态}

一个计数器被修改一百万次，重启时逐条重放这些命令成本很高。AOF重写从当前数据构造能够重建相同状态的较短基础文件，把长修改历史压缩为状态表示。应用如果需要原S与原结果，就让这些记录作为数据存在，被压缩的命令历史不会持续保留每次业务动作的审计含义。

Redis 7以后使用分段AOF结构。基础段BASE保存一次重写得到的状态，增量段INCR承接修改，manifest记录当前应加载的文件及其顺序。Redis 8.10.2继续使用这一结构。基础段可以采用RDB格式的前导数据，形成混合持久化，后面的增量段保留AOF命令流。扩展名和清单描述共同帮助识别文件，恢复者按清单组织完整集合。

\CEighteenFigure{ch18-09}{AOF重写中的新基础段、新增量段与manifest发布}{fig:18-aof-rewrite}

重写开始时，父进程先建立新的增量段，让重写期间发生的写入有明确去处，然后fork子进程生成新的基础段。子进程围绕fork时点的数据写出BASE，父进程继续把后续修改写入新的INCR。子进程成功结束后，父进程把新基础段放到最终名称，构造包含新BASE及所需INCR的新manifest，并持久发布清单。旧文件转为历史文件，清理发生在新清单成功发布之后~\cite{redis8102-source}。

一份简化的文件清单如下：

\begin{BookCode}
file appendonly.aof.3.base.rdb seq 3 type b
file appendonly.aof.3.incr.aof seq 3 type i
file appendonly.aof.4.incr.aof seq 4 type i
\end{BookCode}

这里的\code{seq}是文件序号，\code{type b}表示基础段，\code{type i}表示增量段。当前版本在已知相应位置时，还可以记录复制偏移等附加字段，恢复工具需要识别该版本允许的清单格式。文件序号、复制offset与业务S分别编号不同对象，恢复时分别比较。

假设新BASE已经写完，但新manifest尚未发布，旧manifest及其引用文件是完整加载入口。若新manifest已经可靠发布，恢复按新BASE接上需要的INCR。提前删除旧清单还引用的增量段，会把一个本来能够完成的重写变成缺文件。备份工具因此取得一致的清单与文件集合，或者使用明确的快照、复制及备份协议。

一次重写可以分成三个可恢复的观察点。第一点，当前清单继续以旧BASE为起点，包含此前增量及已纳入清单的新INCR，新BASE尚未完成。第二点，新BASE完成并取得最终名称，清单却还没有选择它，单看文件修改时间选最新BASE会跳过当前清单的选择。第三点，新清单已经持久发布，旧历史文件即使尚未清理，也只是额外占用空间。发布记录把一组文件变成一个可解释的数据版本。

\CEighteenFigure{ch18-23-manifest}{重写的三个观察点分别由已发布清单选择恢复文件集合}{fig:18-manifest-points}

这个组织还解释了AOF重写与业务保留策略的关系。应用主动删除旧S的原结果后，新的BASE反映删除后的数据，历史INCR最终也会被回收。磁盘上暂时残留的旧段由文件生命周期管理，原结果查询则只承诺业务保留规则允许查询的数据。身份回收和准入规则先决定允许忘记什么，重写与文件清理再回收相应表示。

重写减少文件总量，却可能增加瞬时空间：旧BASE、旧INCR、新BASE和重写期间新增的INCR有一段共同存在的时间，内存还承担fork及写时复制。恢复吞吐、磁盘余量和后台任务竞争由此进入容量设计。AOF目录要给新旧文件的安全替换留出空间。

\subsection{启动到底加载哪份数据}

常规启动路径中，启用AOF就按AOF清单加载基础段和有序增量段；没有启用AOF时，使用普通RDB加载路径。两者同时配置时，加载器按这一分支选择入口。采用RDB前导段的BASE本来就是AOF清单的一部分，普通\code{dump.rdb}与它属于不同入口。

Redis 8.10.2还提供显式\code{preload-file}路径。配置指定带类型的\code{rdb:}或\code{aof:}来源时，启动逻辑优先按这个来源预加载，失败会导致启动失败。讨论恢复选择时，先检查有没有显式来源，再检查正常AOF开关与清单，最后确定预期加载的数据~\cite{redis8102-source}。

\CEighteenFigure{ch18-10}{显式预加载与常规RDB、AOF启动分支}{fig:18-startup-choice}

AOF加载按BASE和增量序列重建状态。清单引用的文件缺失，意味着这条恢复历史少了一段，加载器按错误处理。某个更早的BASE即使能够独立解码，也只代表更早状态。恢复者先确认缺失范围和独立证据，再从完整备份、其他副本或冻结对账流程继续，决定何时重新开放支付受理。

\subsection{末尾截断与格式损坏}

追加命令中途退出，最后的INCR可能留下不完整请求。配置项\code{aof-load-truncated}取\code{yes}时，加载器在受支持的末尾截断情形下加载完整前缀，并截去不完整尾部。这里恢复的是最后一个可解析前缀，尾部动作是否已经在独立平台发生，由原C/T协议判断\footnote{EOF表示读到文件结尾。可恢复的尾部截断与命令流中间的格式损坏有不同判断路径，多段AOF还要满足最后文件等条件。}。

格式损坏与读到文件结尾不同。Redis 8.10.2另有\code{aof-load-corrupt-tail-max-size}控制有限尾部格式损坏的恢复，默认值0关闭这一容错。该选项与EOF截断选项分别处理不同条件。多段AOF中，允许的尾部恢复还受到最后文件的限制，非末段损坏会使后续流的解释失去可靠前提~\cite{redis8102-source}。

操作顺序从保留原坏文件、清单和摘要开始，在独立副本上调用检查与修复工具。恢复后核对完整S、原输入、P/C/T、结果与未决责任，并与平台和消息域对账。若本地尾部被截去的是$R_J$与O/EV的共同提交，恢复者依据原平台操作重建同一结果及事件。文件修复工具判断字节流能否加载，X1、履约F501、原交接回执和关闭会话的业务关联由购买系统恢复协议核对。
\FloatBarrier
